\section{Conclusion}

We have described a sparse attention layer whose router is a learned orthogonal rotation acting on a fixed orthonormal codebook, which reads a RoPE-free content subspace, and whose excluded tokens remain visible to the softmax as coarse residue summaries. We gave a non-collapse property of the router's partition, an equidistribution argument for why angular routing should not read rotary dimensions, a torus-confinement result explaining what goes wrong when it does, an invariance result showing that the bucket load belongs to the data rather than to the router, and a complexity and cache analysis that avoids the two overclaims most common in this area --- that expert routing of projections reduces the quadratic term, and that sparse attention reduces KV cache capacity. Neither is true.

The more useful output of this work is probably the four corrections in Section~\ref{sec:studies}, three of which were corrections to our own claims. Routing on rotary dimensions does not degrade toward uniform hashing but toward a capacity ceiling, which is worse and gets relatively worse as the codebook grows. A balancing loss on the router cannot work, which killed an ablation we had already designed. The chunking that the complexity figure depends on is faithful only under a condition the figure omits, which couples routing width to context length. And at initialisation every ablation arm produces an identical pattern, which means the remaining claims cannot be tested cheaply.

We have trained no language model, and we think a proposal is more useful stated that way than dressed in results it does not have. What the released code does provide is the reference implementation, a test suite that checks causality by autograd rather than by inspection, and the studies above. The cheapest decisive experiment is still E0, which requires no training; the revised version in Section~\ref{sec:e0} measures two quantities the original would have missed. We would run it before anything else.

\paragraph{Reproducibility.} The implementation, tests and studies are released at \url{https://github.com/Alessio2405/phasor}. Running \texttt{pytest tests/} executes 128 tests covering orthogonality under optimisation, causality of every ablation arm by autograd, agreement between the training-time and decoding-time summary paths, and the numerical counterparts of each proposition. The following regenerates every table and figure in Section~\ref{sec:studies} from scratch:
\begin{center}\small\ttfamily
python experiments/numerical\_studies.py \&\& python experiments/make\_tables.py
\end{center}

\paragraph{Acknowledgements.} The RoPE-decoupling and residual-fallback ideas in this paper originated in an informal design discussion; the author thanks the interlocutors for pressure-testing them.
